\documentclass[a4paper]{article}

\usepackage{amsmath, amssymb}
\usepackage{multicol}
\setlength{\columnsep}{1cm}
\usepackage[margin=0.75in]{geometry}
\usepackage{graphicx}
\graphicspath{ {../../Images/MATH11218/} }
\usepackage[breaklinks, hidelinks]{hyperref}
\usepackage{tikz}
\usepackage{stix}
\usepackage{ragged2e}
\usepackage{polynom}

\title{MATH11218 \\ Methods and Basically Specialist Revision}
\author{Authored by: Lachlan J. W. Wyer}
\date{\today}

\begin{document}
	\maketitle
	\tableofcontents

	\section{Introduction}
	\subsection{Indices:}
	Within the expression of $a^{m}$, a is considered to be the base with the power m being the index. The reason that this is considered useful is in scientific notation to expression large numbers, or simplify repeated multiplication, fractions, roots or polynomials. E.g. $2\times 2\times 2\times 2\times 2\times 2=2^{6}$. With these there are a number of different index laws that goes with it:
	\begin{enumerate}
		\item $a^{m}\cdot a^{n}\;=\;a^{m+n}$
		\item $\frac{a^{m}}{a^{n}}\;=\;a^{m-n}$
		\item $(a^{m})^{n}\;=\;a^{m\cdot n}$
		\item $(ab)^{m}\;=\;a^{m}b^{m}$
		\item $(\frac{a}{b})^{m}\;=\;\frac{a^{m}}{b^{m}}$
		\item $a^{\frac{m}{n}}\;=\;\sqrt[n]{a^{m}}$
		\item $a^{0}\;=\;1$
	\end{enumerate}

	\subsection{Poynomials:}
	A polynomial explains/denotes the degree or curviness that a line is at. Such as a constant can be described as being degree 0, linear degree 1, quadratic degree 2, cubic degree 3 and  so on.
	\subsubsection*{Quadratics}
	With quadratics it is important to understand that it can contain multiple different types of roots. For example, iff: $\Delta=b^{2}-4ac$
	\begin{itemize}
		\item $\Delta\;>\;0$, when it contains 2 real roots, e.g. $(x+a)(x+b)$.
		\item $\Delta\;=\;0$, when there is 1 real root, e.g. $(x+a)^{2}$.
		\item $\Delta\;<\;0$, when no real roots are found, e.g $(x+aj)(x+bj)$.
	\end{itemize}
	\subsubsection*{Polynomials with Degrees Greater than 2}
	When a polynomial contains a degree that is greater than 2, it roots can be found through a process called polynomial long division.
	\medskip
	\begin{center}
		\polylongdiv{x^{3}+3x^{2}-10x-24}{x+2}\\
		\smallskip
		\begin{align}
			\therefore\;\;x^{3}+3x^{2}-10x-24\;&=\;(x+2)(x^{2}-x-12)\\
			&=\;(x+2)(x-4)(x+3)\\
			x\;&=\;-3,\;-2,\;4
		\end{align}
	\end{center}

	\subsection{Algebraic Functions:}
	Within fractional functions there are two main types:
	\begin{enumerate}
		\item Proper fractions, where the denominator is smaller then the numerator.
		\item Improper fractions, where the denominator is less or equal to the numerator.
	\end{enumerate}
	It is important to note, that the greater, equal to, and less than are only relative to polynomial degree. Some parts to with division of fraction is just the multiplication of the recipricol, $\frac{a}{b}\;\div\;\frac{c}{d}\;=\frac{a}{b}\;\times\;\frac{d}{c}$.
	\begin{align}
		&\frac{2}{x+1}\;+\;\frac{4}{x+2}\\
		&=\;\frac{2(x+2)\;+\;4(x+1)}{(x+1)(x+2)}\\
		&=\;\frac{2x+4+4x+4}{(x+1)(x+2)}\\
		&=\;\frac{6x+8}{(x+1)(x+2)}\\
		\therefore\;\;x&\not =\;-2,\;-1
	\end{align}
\end{document}
